\section{First-campaign evaluation and implementation comparison}\label{sec:benchcomparison}
\subsection{Evidence cutoff and measurement boundaries}
The campaign recorded signing and verification latencies, initialization costs, and LULA I fixed-snapshot behavior. Its final retained record was frozen at 14:52 UTC on 12 October 2026. All four controllers have completion markers, and all 81 recorded processes completed with exit code zero. We independently checked the 477 retained file digests and recomputed distributions from 64 campaign CSV files, excluding smoke tests. Process completion establishes that scheduled jobs ended; it does not establish conformance or completion of the original experimental protocol. This first campaign lacked quiet-window trials, equivalent backends and a common-security cohort. The later workload gate and two-key prepared measurements are reported separately in \Cref{sec:stability}. The tables use raw observations subject to the inclusion criteria stated here, rather than relying on the campaign's summary labels.

The physical host was an Intel Core Ultra 7 255U with 32 GiB RAM running Windows 11 Pro, build 26300, under the Balanced power plan. Each measured campaign child job was pinned to logical processor 10 (affinity mask 1024), at BelowNormal priority, with a 2.75 GiB committed-memory ceiling. Its SMT sibling was not isolated. Other user workloads continued. The controller sampled system load before and after jobs, rather than continuously. CPU frequency and temperature were not controlled. Thus the experiment fixes a processor affinity but does not meet a quiet, isolated-core protocol.

LULA I used Rust 1.97.1 and \texttt{sha2} 0.10.9, with optimization level 3, one codegen unit and link-time optimization. The C adapters used Clang 22.1.8 with \texttt{-O3 -mno-avx}; OpenSSL 3.5.7 was supplied by the Git for Windows DLL and retained its runtime dispatch. Equivalent source-level flags do not make these hash backends equivalent, and instrumentation also differs. We report implementation observations; we do not attribute between-scheme timing differences solely to cryptographic design.

Elapsed nanoseconds were measured with high-resolution platform counters. Signing and verification were timed separately around individual API calls. Message construction and process launch are outside these intervals. The OpenSSL adapter initializes its API contexts before timing; internal hash-context initialization performed by a signing call remains inside that call's interval. The adapter supplies no deterministic-mode parameter, so OpenSSL uses its default randomized (hedged) signing mode~\cite{openssl}; LULA I signing is deterministic. OpenSSL uses the whole-message SLH-DSA interface with empty context~\cite{openssl}; the LULA I interface hashes the complete input internally. A 32-byte input is a raw digest-sized message, not a standardized prehash-mode measurement. No separate external-prehash cost was measured. LULA I uses replicates 0 through $N-1$, whereas the OpenSSL driver uses 1 through $N$; these are equal-length cells, not paired-message trials. LULA I uses one key throughout this campaign; OpenSSL generates one key per message-length cell.

The median is the midpoint of the two central order statistics when $N$ is even. Percentiles use nearest rank, $x_{\lceil pN\rceil}$ in one-based sorted order. This differs from the earlier verification replay's nearest-index convention. Untimed negative-control sentinels are excluded. We retain observed maxima and sample counts without interpreting them as mathematical runtime bounds.

\subsection{Prepared signing and whole-message verification}
\begin{table}[t]
\caption{First-campaign signing (S) and verification (V) elapsed latencies, in milliseconds. Signing-query reference points differ; see \Cref{sec:operating}. Max is the maximum observed, not a worst-case bound.}\label{tab:comparison-latencies}
\centering\small\setlength{\tabcolsep}{3pt}
\begin{tabular}{llrrrrrr}\toprule
Cell & Op. & $N$ & Min & Median & p95 & p99 & Max\\\midrule
L/32 B & S & 1000 & 4.151 & 4.600 & 5.332 & 5.718 & 30.530\\
L/32 B & V & 1000 & 0.129 & 0.141 & 0.166 & 0.190 & 5.629\\
L/1 KiB & S & 1000 & 3.676 & 3.920 & 4.242 & 4.843 & 12.162\\
L/1 KiB & V & 1000 & 0.125 & 0.131 & 0.142 & 0.179 & 0.279\\
L/64 KiB & S & 1000 & 3.685 & 3.968 & 4.696 & 5.249 & 10.778\\
L/64 KiB & V & 1000 & 0.151 & 0.161 & 0.189 & 0.228 & 0.388\\
F/32 B & S & 1000 & 10.522 & 12.004 & 14.181 & 21.761 & 44.068\\
F/32 B & V & 1000 & 0.615 & 0.719 & 0.847 & 1.024 & 3.913\\
F/1 KiB & S & 1000 & 12.011 & 13.925 & 18.339 & 25.634 & 36.288\\
F/1 KiB & V & 1000 & 0.684 & 0.823 & 1.018 & 1.549 & 2.744\\
F/64 KiB & S & 1000 & 10.642 & 11.498 & 13.626 & 16.378 & 383.005\\
F/64 KiB & V & 1000 & 0.659 & 0.741 & 0.897 & 1.084 & 17.350\\
S/32 B & S & 100 & 251.915 & 268.616 & 330.976 & 367.905 & 408.373\\
S/32 B & V & 100 & 0.234 & 0.268 & 0.321 & 0.362 & 0.500\\
S/1 KiB & S & 100 & 249.424 & 259.794 & 327.080 & 362.245 & 373.314\\
S/1 KiB & V & 100 & 0.228 & 0.262 & 0.311 & 0.380 & 0.429\\
S/64 KiB & S & 100 & 232.917 & 247.069 & 259.209 & 267.189 & 279.123\\
S/64 KiB & V & 100 & 0.290 & 0.319 & 0.358 & 0.436 & 0.446\\
\bottomrule\end{tabular}
\smallskip\par\raggedright\footnotesize L: prepared LULA I; F/S: OpenSSL SLH-DSA-SHA2-128f/128s. Message sizes refer to raw bytes supplied to the signing API. Values rounded to 0.001 ms; raw nanoseconds are retained.
\end{table}
\Cref{tab:comparison-latencies} summarizes the admitted latency samples. The prepared LULA I signing medians are 4.600, 3.920 and 3.968 ms for 32-byte, 1 KiB and 64 KiB messages; verification medians are 0.141, 0.131 and 0.161 ms. The smaller 1 KiB signing median than the 32-byte median is an observation from different execution intervals, not negative message-hashing cost. System load before/after those cells was 32/6\%, 41/16\% and 8/8\%, respectively. Endpoint load samples cannot explain individual outliers.

The signature lengths returned by the adapters are 5,720 bytes for LULA I, 17,088 for SHA2-128f and 7,856 for SHA2-128s. LULA I's public key is 32 bytes. Both SLH-DSA public keys are 32 bytes by the standard format~\cite[Section~9.1 and Table~2]{fips205}; the OpenSSL adapter did not export their serialized keys. These values establish an observed size/latency tradeoff for the named implementations, but not comparable unforgeability bounds. LULA I uses a conditional $2^{40}$-query objective; SHA2-128s/128f use the published $2^{64}$-query targets (\Cref{tab:security-model}).

\begin{table}[t]
\caption{Signature components in bytes. SLH-DSA decomposition follows its standardized parameters; LULA I follows the specified wire encoding. All values sum to the observed signature lengths. Query points: LULA I conditional $2^{40}$; SLH-DSA $2^{64}$ (\Cref{tab:security-model}).}\label{tab:comparison-bytes}
\centering\small
\begin{tabular}{lrrr}\toprule
Component & LULA I & SHA2-128s & SHA2-128f\\\midrule
Randomizer &16&16&16\\
FORS-family evidence &2,720&2,912&3,696\\
Hypertree chain values &2,304&3,920&12,320\\
Hypertree authentication &656&1,008&1,056\\
Explicit search counters &24&0&0\\\midrule
Total &5,720&7,856&17,088\\\bottomrule
\end{tabular}
\end{table}
\Cref{tab:comparison-bytes} separates compressed-component and hypertree costs. LULA I transmits 17 FORS paths, four lower 32-chain WOTS vectors and one upper 16-chain vector. SHA2-128s and SHA2-128f each use 35-chain WOTS vectors, at seven and 22 layers respectively. Removing chain values changes the construction; retaining precomputed nodes changes signing work without shortening this wire encoding. For valid LULA I verification the instrumented compression counts are 3,265, 3,280 and 4,288 at the three message lengths. Context-prefix initialization executes one additional compression outside that counter. Hardware CPU cycles and corresponding counters for the other adapters are unavailable.

\Cref{fig:comparison-distributions} shows the retained signing and verification distributions, including their observed tails.
\begin{figure}[!htbp]
\centering\includegraphics[width=\textwidth]{figures/comparison-distributions.pdf}
\caption{Empirical cumulative distributions of separate signing and verification calls on 32-byte messages. LULA I has a prepared cache. Query points are conditional $2^{40}$ for LULA I and published $2^{64}$ for SLH-DSA (\Cref{tab:security-model}). Curves describe retained samples under varying host load and differing backends; they do not establish a security-normalized ranking.}\label{fig:comparison-distributions}
\end{figure}

\subsection{Cache preparation, restart and complete cache loss}
One affinity-constrained LULA I key-generation call, including complete cache preparation, took 2,216.118 seconds and recorded 45,198,621,044 instrumented compressions. The harness then spent 23.908 seconds writing its persisted fixture. The controller's whole-process wall time was 2,240.179 seconds, including startup and exit effects; it is not an additional setup phase. This one-worker observation is separate from the historical multiworker runs.

Cache payload is 75,497,456 bytes. The experimental fixture occupies 75,497,573 bytes, including 117 bytes of framing and key material. During preparation, the recorded whole-process peak working set was 122,789,888 bytes and private memory 143,585,280 bytes. These are process metrics for that run, not per-operation memory or additional transmitted data. A phase-specific, algorithm-normalized memory cohort remains missing.

\begin{table}[t]
\caption{LULA I recovery observations on 32-byte messages, in seconds. The two cache-loss calls share one process; each is shown individually. Paired sums include fixture loading and signing, but exclude OS process launch.}\label{tab:comparison-recovery}
\centering\small
\begin{tabular}{P{57mm}rrr}\toprule
Interval & $N$ & Median & Observed range\\\midrule
Fresh process: fixture load &10&6.185605&5.906147--8.333199\\
Fresh process: first Sign after load &10&0.003959&0.003846--0.005281\\
Paired load plus first Sign &10&6.189459&5.910034--8.338480\\
Cache absent: first Sign &1&\multicolumn{2}{c}{346.189046}\\
Cache still absent: second Sign &1&\multicolumn{2}{c}{330.789418}\\\bottomrule
\end{tabular}
\end{table}
Ten fresh processes loaded the same honest full-cache fixture and signed the same first message. \Cref{tab:comparison-recovery} reports the median across the per-process load-plus-sign sums; adding marginal medians would be incorrect. The loader checks dimensions and reads the structure, but does not authenticate every cached node. Persistence is an experimental harness codec, rather than a native optimized restore API. These timings concern honest saved bytes, consistent with the threat model.

In a separate process the cache-loss driver loaded only seed/root material, in 0.235 ms, and installed no layer cache. The first signature took 346.189 seconds and recorded 8,602,614,321 instrumented compressions. The next, on another message, took 330.789 seconds and recorded 8,602,614,270. The signer did not automatically retain a replacement top cache. Two calls in one process are not two independent recovery trials; neither their minimum nor their median is the first-signature latency. This new fixture and execution context differ from the historical 703.996-second observation, so their difference is not a controlled speedup.

The independent Python verifier accepted all 14 retained LULA I vectors and rejected changed messages and signatures. They comprise three prepared-cache input lengths, ten first signatures after fixture reload, and the first cache-loss signature. For the same 32-byte message and key, 12 signatures out of those 14 campaign vectors---prepared, reloaded and absent-cache instances---have identical complete signature bytes. The second cache-loss signature passed native verification but has no separately retained output vector. The completed branch-schedule experiment is described below. General rollback transparency follows from the theorem's hypotheses; these selected executions test its implementation. \Cref{fig:comparison-recovery} displays the separation between prepared signing, fixture recovery and the two consecutive cache-loss calls.

\begin{figure}[!htbp]
\centering\includegraphics[width=\textwidth]{figures/comparison-recovery.pdf}
\caption{LULA I signing and recovery intervals on a logarithmic elapsed-time axis. Bars show observed ranges where multiple samples exist. Cache-loss points are consecutive individual calls. Loading and first signing are paired per process. No intermediate cache sizes were measured, so these points do not define a cache-memory/latency curve.}\label{fig:comparison-recovery}
\end{figure}

OpenSSL key generation was measured once per message-length cell. SHA2-128f took 4.121, 5.666 and 4.578 ms, and SHA2-128s took 43.193, 36.954 and 51.577 ms, in 32-byte/1 KiB/64 KiB order. These isolated API intervals exclude provider loading and context creation. Comparing them directly with LULA I's key generation plus full preprocessing would conflate distinct setup boundaries.

\subsection{SHRINCS implementations and restoration}
\begin{table}[t]
\caption{Version-specific upstream C SHRINCS B32 observations, in milliseconds. Stateful (SF) calls advance the index; fallback (FB) is explicitly forced under the same key. These are descriptive implementation results.}\label{tab:comparison-b32}
\centering\small\setlength{\tabcolsep}{3pt}
\begin{tabular}{llrrrrrr}\toprule
Cell & Op. & $N$ & Min & Median & p95 & p99 & Max\\\midrule
SF/32 B & S & 200 & 45.595 & 48.520 & 59.049 & 71.829 & 96.670\\
SF/32 B & V & 200 & 0.106 & 0.120 & 0.142 & 0.273 & 0.642\\
SF/1 KiB & S & 200 & 45.810 & 47.934 & 52.885 & 61.069 & 145.431\\
SF/1 KiB & V & 200 & 0.106 & 0.117 & 0.136 & 0.216 & 0.407\\
SF/64 KiB & S & 200 & 45.632 & 47.244 & 55.264 & 102.388 & 125.685\\
SF/64 KiB & V & 200 & 0.140 & 0.153 & 0.163 & 0.233 & 0.270\\
FB/32 B & S & 20 & 423.866 & 471.961 & 553.538 & 616.742 & 616.742\\
FB/32 B & V & 20 & 0.444 & 0.448 & 0.493 & 0.537 & 0.537\\
FB/1 KiB & S & 20 & 431.423 & 497.561 & 676.710 & 799.294 & 799.294\\
FB/1 KiB & V & 20 & 0.445 & 0.447 & 0.454 & 0.511 & 0.511\\
FB/64 KiB & S & 20 & 432.752 & 3558.347 & 9576.367 & 10182.099 & 10182.099\\
FB/64 KiB & V & 20 & 0.489 & 0.497 & 0.547 & 0.565 & 0.565\\
\bottomrule\end{tabular}
\smallskip\par\raggedright\footnotesize SF: 200 successive signatures, 324--3,508 bytes observed. FB: 3,680 bytes. Small fallback samples make p99 equal to the observed maximum. Parameters, backend and security obligations differ from the primary table.
\end{table}
\Cref{tab:comparison-b32} describes the upstream C B32 implementation at commit \texttt{345b0a64}~\cite{shrincsc}. Its node size is 16 bytes, with WOTS+C width 256 and 16 chains, stateful height 210, a total stateless hypertree height of 32 across four height-8 layers, and a height-17 PORS+FP tree revealing 11 selected leaves. Its public key is 32 bytes; observed stateful signatures span 324--3,508 bytes over 200 successive calls, and forced fallback signatures are 3,680 bytes. It is distinct from the pinned SHRINCS draft at commit \texttt{a5891ec5}, whose public key is 48 bytes and whose stateful and stateless encodings are 548--4,619 and 5,777 bytes~\cite{shrincsdraft}. We do not substitute one version's parameters or proof obligations for another's.

The B32 fallback cells explicitly invoke stateless signing under the original key. They are not detected-loss measurements. Ten fresh-process stateful signatures under the same key had a signing median of 46.243 ms, excluding fixture loading. The detected-loss adapter passed an incorrect seed layout to its restore routine and replaced the public key, as confirmed by the saved vectors; those observations are excluded from same-key recovery comparison. Three completed campaign processes also exercised a scripted transition near the stateful endpoint. The driver installs stored index 207, takes four stateful calls, and then selects fallback at stored index 211. The emitted stateful indices are 208 through 211, with lengths 3,636, 3,652, 3,668 and 3,668 bytes. This is a finite transition exercise, not a naturally exhausted signing history. The driver does not call the stateful API at its switch point to observe a refusal; no refusal latency or certified 210-signature budget follows.

The silent-rollback adapter reset the state before signing distinct messages and reused a stateful index. Recording valid signatures after reuse does not show unforgeability across branches.

Silent restoration of a counter-based signer can return to a previously used one-time address. Distinct accepted target-sum vectors at that address can expand the forward-reachable set (\Cref{sec:security}); the size of that set depends on the exposures, so reuse is not assigned a universal forgery cost here. Statelessness avoids this consumption-index rewind; restoration of signing randomness still needs a joint transcript argument. For LULA I, the claim is specifically honest-cache transparency for successfully generated keys and preservation of an applicable reference bound, not a completed proof of that bound. No forgery was run against a real SHRINCS implementation in this work, including its earlier campaigns. The increased forgery exposure from a reused counter in target-sum WOTS+C is a model argument, not an experiment against SHRINCS. Reliable loss detection and silent restoration therefore remain different cases, regardless of their elapsed times.

The third-party Go SHRINCS adapter~\cite{hashsiggo} has completed cells, but its parser accepted appended bytes and one tested counter-header mutation. Those controls do not establish a fresh-message forgery; they do prevent calling the adapter fully validated. Its timings are retained as diagnostic data rather than primary comparison rows. The currently inspected code hashes the complete message; no claim of a message-prefix defect is inferred from a different historical implementation.

\subsection{Compact reconstructions and auxiliary adapter measurements}
\begin{table}[t]
\caption{Auxiliary compact-adapter signing (S) and verification (V) elapsed latencies, in milliseconds. These complete timing cells are not measurements of the authors' optimized implementations.}\label{tab:comparison-kn}
\centering\small\setlength{\tabcolsep}{3pt}
\begin{tabular}{llrrrrrr}\toprule
Cell & Op. & $N$ & Min & Median & p95 & p99 & Max\\\midrule
R/32 B & S & 20 & 2635.074 & 3613.611 & 4098.247 & 4128.473 & 4128.473\\
R/32 B & V & 20 & 0.451 & 0.542 & 1.008 & 1.118 & 1.118\\
R/1 KiB & S & 20 & 2086.219 & 3241.905 & 4856.846 & 5543.975 & 5543.975\\
R/1 KiB & V & 20 & 0.482 & 0.712 & 1.259 & 1.260 & 1.260\\
R/64 KiB & S & 12 & 2051.909 & 3381.547 & 5127.710 & 5127.710 & 5127.710\\
R/64 KiB & V & 12 & 0.514 & 0.670 & 1.248 & 1.248 & 1.248\\
G/32 B & S & 12 & 1084.598 & 1897.784 & 4450.316 & 4450.316 & 4450.316\\
G/32 B & V & 12 & 0.883 & 0.905 & 1.781 & 1.781 & 1.781\\
G/1 KiB & S & 12 & 1168.415 & 2061.765 & 3000.875 & 3000.875 & 3000.875\\
G/1 KiB & V & 12 & 0.828 & 0.910 & 1.656 & 1.656 & 1.656\\
G/64 KiB & S & 12 & 1754.491 & 12121.972 & 78265.188 & 78265.188 & 78265.188\\
G/64 KiB & V & 12 & 0.986 & 1.310 & 11.912 & 11.912 & 11.912\\
\bottomrule\end{tabular}
\smallskip\par\raggedright\footnotesize R: 3,952-byte reconstruction; G: 4,036-byte third-party Go adapter. Small N limits tail interpretation; p99 equals the observed maximum for these cells. No security-normalized comparison is inferred.
\end{table}
The compact configurations discussed by Kudinov and Nick~\cite{kn} motivate two additional measured adapters, distinguished in \Cref{tab:comparison-kn}. The 3,952-byte variant is our noninteroperable reconstruction, not the authors' optimized implementation: it uses height 44 across four layers, eight plain FORS trees of height 16, 16 width-256 WOTS+C chains per layer, and private hash domains. The 4,036-byte variant is third-party Go code~\cite{hashsiggo}, with total height 40 across five layers and PORS+FP selecting 11 leaves. Its PORS exponent $a=14$ gives $11\cdot2^{14}$ leaves in a height-18 tree; $a$ is not that tree's height. Neither row has a checked reduction and conformance argument for its exact implementation. Their public keys are 32 bytes in the exported vectors.

Digest-sized signing medians are 3,613.611 ms for the reconstruction and 1,897.784 ms for the Go adapter; their verification medians are 0.542 and 0.905 ms. The prior efficiency-core smoke and the later logical-processor-10 campaign differ in core assignment and host activity. Their difference does not identify a backend or architectural cause. Even an observed maximum supplies no guaranteed upper bound on this reconstruction or the authors' code.

The Go adapter's ``prepared'', ``fresh'' and ``missing'' labels all invoke the same uncached algorithm after loading key constants. They are not interventions on different cache states. The reconstruction's fresh-process measurements likewise load key material without creating LULA I's persisted-cache mechanism. These adapters' new complete timing cells supply auxiliary implementation observations, without converting the labels into recovery guarantees.

\begin{table}[t]
\caption{Fresh-process sample medians on 32-byte messages, in milliseconds. Loading plus first Sign is paired per process before aggregation; OS process launch is excluded. Initialization contents differ.}\label{tab:comparison-fresh}
\centering\small
\begin{tabular}{lrrrr}\toprule
Loaded fixture & $N$ & Load & First Sign & Paired sum\\\midrule
LULA I, full cache & 10 & 6185.605 & 3.959 & 6189.459\\
C B32, key/state & 10 & 0.453 & 46.243 & 46.682\\
KN-R, key material & 3 & 0.614 & 5577.513 & 5578.127\\
KN-G, key material & 3 & 1.012 & 2011.245 & 2012.257\\
\bottomrule\end{tabular}
\end{table}
\Cref{tab:comparison-fresh} exposes the differing initialization boundaries. Each row uses one saved key in multiple new processes. We calculate each load-plus-sign sum within a process, then take the median across those sums. These are not independent-key trials. LULA I reads its large cache fixture, whereas the auxiliary adapters load key material or key/state structures. The table therefore compares actual first-use observations, rather than treating millisecond signing after a seconds-long load as restart latency.

\begin{table}[t]
\caption{Whole-process peak-memory observations for the 32-byte cells, in MiB. These include setup and the complete call sequence; they are not phase-specific or independent-key distributions.}\label{tab:comparison-memory}
\centering\small
\begin{tabular}{lrr}\toprule
Implementation cell & Working set & Private\\\midrule
LULA I, prepared & 76.645 & 73.203\\
SLH-DSA SHA2-128f & 9.309 & 2.094\\
SLH-DSA SHA2-128s & 9.293 & 2.090\\
C B32, stateful & 6.980 & 1.371\\
C B32, forced fallback & 6.988 & 1.367\\
KN-R reconstruction & 7.973 & 2.383\\
KN-G third-party & 53.746 & 89.824\\
\bottomrule\end{tabular}
\end{table}
\Cref{tab:comparison-memory} reports controller-observed process peaks for the 32-byte cells, including fixture loading, runtime libraries and the full sequence of timed calls. One process per cell supplies each peak. Neither working set nor private memory is per-signature storage or cache payload. The controller combines polled private usage with the available post-exit peak metric. Runtime, allocator and shared-library differences prevent isolating an algorithm-only memory cost from these numbers.

\begin{table}[t]
\caption{Exact signature-component bytes of the measured auxiliary encodings. B32 FB is upstream C fallback; KN-R is our reconstruction; KN-G is the third-party Go adapter. These decompositions are derived from the implementation and sum to observed totals; they do not establish interoperability.}\label{tab:comparison-aux-bytes}
\centering\small
\begin{tabular}{lrrr}\toprule
Component & B32 FB & KN-R & KN-G\\\midrule
Sibling component root &16&0&0\\
Randomizers &160&32&32\\
Few-time revealed secrets &176&128&176\\
Few-time authentication &1,776&2,048&1,888\\
Hypertree chain values &1,024&1,024&1,280\\
Hypertree authentication &512&704&640\\
Layer search counters &16&16&20\\\midrule
Total &3,680&3,952&4,036\\\bottomrule
\end{tabular}
\end{table}
In \Cref{tab:comparison-aux-bytes}, PORS authentication allocations are fixed and padded in B32 FB and KN-G; KN-R transmits eight complete height-16 FORS paths. B32 stateful evidence instead contains a 16-byte sibling root, 32-byte randomizer, four-byte counter, 256 chain-value bytes and $16\min(q,210)$ authentication bytes for the emitted index $q$. This yields 324 bytes at $q=1$, 3,508 at $q=200$, and 3,668 at either endpoint $q=210$ or $211$. Encoding length and a source endpoint are not an independently established safe signing budget.

\subsection{Branch-schedule observations}
\begin{table}[t]
\caption{LULA I fixed-snapshot branch signing latencies, in milliseconds. Verification succeeds for all driver outputs; no separate verification durations were retained.}\label{tab:comparison-branches}
\centering\small\setlength{\tabcolsep}{3pt}
\begin{tabular}{lrrrrrr}\toprule
Branch & $N$ & Min & Median & p95 & p99 & Max\\\midrule
32 B/A & 100 & 4.358 & 6.536 & 9.273 & 9.754 & 9.971\\
32 B/B & 200 & 4.900 & 8.172 & 13.581 & 29.668 & 30.590\\
1 KiB/A & 50 & 4.034 & 6.175 & 8.706 & 9.116 & 9.116\\
1 KiB/B & 100 & 4.995 & 8.203 & 18.714 & 36.021 & 51.109\\
\bottomrule\end{tabular}
\end{table}
Two processes per input length loaded the original full-cache fixture or its byte-identical saved copy. Branch A signed message identifiers 0 through 99 and branch B identifiers 0 through 199 for 32-byte inputs. The 1 KiB schedules used 0 through 49 and 0 through 99. Independent comparison of complete serialized signatures found equality on all 100 and 50 overlapping messages; branch B also produced 100 and 50 additional signatures. All signatures are 5,720 bytes and the driver verified each.

The saved and live cache files have the same recorded SHA-256 digest. Signing reads an immutable cache reference and does not advance an authorization index, so these are separate evaluations from one honest fixed fixture. They are not measurements of a progressively mutated cache, an aged-cache transition, or an adaptive forgery attempt. The branch sign latencies in \Cref{tab:comparison-branches} are recomputed from retained nanoseconds using the ordinary midpoint median. The campaign's separate branch summary used the upper middle value for even sample counts. These finite observations support byte invariance in the chosen schedules; rollback unforgeability remains conditional on the common reference-signer bound.

\subsection{Supported tradeoff and unresolved comparisons}
\begin{table}[!htbp]
\caption{Eight-candidate overview. Every measured latency is a 32-byte-message median from the first campaign only; the gated follow-up is separate. Signature and public-key sizes are bytes; S/V latencies are milliseconds. M: measured; D: derived; S: source specification or documented availability; P: unestablished premise. n.m.: not measured. Signing-query references are distinct targets or conditional scenarios, not certified capacities of the auxiliary adapters.}\label{tab:comparison-overview}
\centering\scriptsize\setlength{\tabcolsep}{2pt}
\begin{tabular}{P{20mm}P{25mm}rrrrr}\toprule
Query reference & Construction & $|\sigma|$ & $|PK|$ & Sign & Verify & $N_{S,V}$\\\midrule
$2^{24}$ target\textsuperscript{S} & Draft SHA2-128-24\textsuperscript{S} & 3,856\textsuperscript{S} & 32\textsuperscript{S} & n.m.\textsuperscript{S} & n.m.\textsuperscript{S} & n.m.\textsuperscript{S}\\
conditional $2^{40}$\textsuperscript{S} & LULA I prepared\textsuperscript{S}\newline bound open\textsuperscript{P} & 5,720\textsuperscript{M} & 32\textsuperscript{M} & 4.600\textsuperscript{M} & 0.141\textsuperscript{M} & 1000\textsuperscript{M}\\
$2^{64}$ target\textsuperscript{S} & SLH-DSA SHA2-128s\textsuperscript{S} & 7,856\textsuperscript{M} & 32\textsuperscript{S} & 268.616\textsuperscript{M} & 0.268\textsuperscript{M} & 100\textsuperscript{M}\\
$2^{64}$ target\textsuperscript{S} & SLH-DSA SHA2-128f\textsuperscript{S} & 17,088\textsuperscript{M} & 32\textsuperscript{S} & 12.004\textsuperscript{M} & 0.719\textsuperscript{M} & 1000\textsuperscript{M}\\
Not certified\textsuperscript{P} & SHRINCS C B32 stateful\textsuperscript{S} & 324--3,508\textsuperscript{M} & 32\textsuperscript{M} & 48.520\textsuperscript{M} & 0.120\textsuperscript{M} & 200\textsuperscript{M}\\
Not certified\textsuperscript{P} & SHRINCS C B32 fallback\textsuperscript{S} & 3,680\textsuperscript{M} & 32\textsuperscript{M} & 471.961\textsuperscript{M} & 0.448\textsuperscript{M} & 20\textsuperscript{M}\\
$2^{40}$ example\textsuperscript{S} & KN components: C reconstruction\textsuperscript{S}\newline adapter bound open\textsuperscript{P} & 3,952\textsuperscript{M} & 32\textsuperscript{M} & 3613.611\textsuperscript{M} & 0.542\textsuperscript{M} & 20\textsuperscript{M}\\
$2^{40}$ example\textsuperscript{S} & KN components: third-party Go\textsuperscript{S}\newline adapter bound open\textsuperscript{P} & 4,036\textsuperscript{M} & 32\textsuperscript{M} & 1897.784\textsuperscript{M} & 0.905\textsuperscript{M} & 12\textsuperscript{M}\\
\bottomrule\end{tabular}
\smallskip\par\raggedright\footnotesize Draft row: SP 800-230 initial public draft, Table~1 and Section~3~\cite{sp800230}; SLH-DSA targets and key sizes: FIPS~205~\cite{fips205}. KN query references name the published examples~\cite[Table~1]{kn}, not proven capacities of these implementations. C B32 uses the pinned upstream version~\cite{shrincsc}; its stateful size range covers the 200 successive timed indices, and fallback is explicitly forced. OpenSSL uses hedged signing, LULA I deterministic signing; backends and message schedules differ. The C reconstruction and Go adapter~\cite{hashsiggo} are auxiliary measurements, not the authors' optimized code. Each $N$ counts timed calls per operation, not independent keys. LULA I's lifecycle costs are paired with this row in \Cref{tab:comparison-recovery}.
\end{table}


\noindent The prepared LULA I row must be read with its first-campaign lifecycle costs: 2,216.118 seconds for key generation plus complete cache preparation, 23.908 seconds for fixture persistence, and a median paired loading-plus-first-signing time of 6.189459 seconds over ten fresh processes. The two absent-cache calls took 346.189046 and 330.789418 seconds in one process, without automatic cache rebuilding (\Cref{tab:comparison-recovery}). These are measured first-campaign intervals, not sums across evidence classes. The follow-up's fifteen fresh processes and three individual cold calls remain separate in \Cref{tab:stability-recovery}.

The named implementations' signing-query reference points differ (\Cref{sec:operating}); the LULA I component model at $2^{40}$ queries is not a full security bound, whereas FIPS 205 specifies the SLH-DSA sets for up to $2^{64}$ signatures. LULA I uses approximately 72 MiB of cache payload and costly preparation; cache loss returns signing to hundreds of seconds. This intervention preserves the selected output bytes while changing work by billions of compressions, directly supporting the mechanism explanation.

The evidence does not identify the strongest construction at a common security level. LULA I's reference bound is still a premise; other versions use different algorithms, parameters, parsers and state assumptions. We therefore defer a security-normalized Pareto frontier. Intermediate cache-size experiments, controlled equivalent backends, hardware cycle counts, independently timed prehashing, a broad multi-key cohort, controlled same-key SHRINCS state-loss recovery and a naturally exhausted signing history remain missing. The gated follow-up adds continuous workload monitoring and two LULA I fixture keys, without providing equivalent backends or physical-core isolation (\Cref{sec:stability}). Published latency maxima cannot replace any of these obligations. The supported contribution remains canonical cache transparency and its implementation-level resource tradeoff, rather than a new general principle of hash-based security.

The locally prepared reproducibility bundle contains our source code, retained campaign records, public verification vectors, analysis scripts and models. Our code is licensed under Apache-2.0; our manuscript text, figures and data are licensed under CC BY 4.0. Third-party material retains its own terms and is excluded where no redistribution license was verified. The SHA-256 digest of its code-and-data \texttt{EVIDENCE-MANIFEST.sha256} is
\begin{center}\small\ttfamily
23d0c9a2ecdaf1867876232f2cee4d44\allowbreak
9231d84d111120c55e1b528cd96fa6ec
\end{center}

Public deposit: \texttt{TODO(author): public URL / DOI}. % TODO(author,Q3): publish the prepared licensed bundle and supply its public URL or DOI; repository location is supplied but remains private.
The \texttt{analyze.py} and \texttt{verify\_vectors.py} entry points reaggregate retained measurements and check saved signatures without signing new messages. Separate scripts regenerate the figures and analytical models. Seven historical rows are recoverable only from surviving aggregate records, without per-call reaggregation. Timing reruns require separately pinned upstream dependencies, executable identities and a new output directory. The excluded SHRINCS C dependency prevents a complete C/KN rebuild from this bundle alone. Digests do not reconstruct missing binaries or certify historical compiler inputs.
